\labsection{8}{VLANs: one switch, several networks}{sec:lab08-vlan}

\begin{labproject}{Access, trunk, and hybrid ports --- and what each does to a tag}
\textbf{Coverage.} Chapter~\ref{ch:vlan}.\\
\textbf{Goal.} Separate broadcast domains on shared hardware, carry several VLANs over one link, and predict tagging behaviour before configuring it.

\medskip\noindent\textbf{Part A --- What a VLAN actually separates.}\par
A switch floods broadcasts to every port. A VLAN restricts that flooding to a subset, which means one physical switch behaves as several independent switches.

\begin{keyidea}
Devices in different VLANs cannot reach each other \emph{at all} without a router, regardless of IP addressing. Putting them in the same subnet changes nothing --- the frames never leave their broadcast domain.

That is the point. VLAN separation is enforced at layer 2, below IP, so it cannot be bypassed by reconfiguring a host.
\end{keyidea}

\medskip\noindent\textbf{Part B --- Three port types.}\par

\begin{mltable}
\begin{tabularx}{\linewidth}{@{}>{\bfseries}p{0.13\linewidth}X X@{}}
\toprule
\rowcolor{MLPaleBlue!92}
Port type & Incoming untagged frame & Outgoing frame \\
\midrule
Access & Tagged with the port's PVID & Tag removed \\
Trunk & Tagged with PVID & Tagged, unless the VLAN is the native VLAN \\
Hybrid & Tagged with PVID & Tagged \emph{or} untagged, configured per VLAN \\
\bottomrule
\end{tabularx}
\end{mltable}

\begin{keyidea}
Every port has a \term{PVID} --- the VLAN assigned to arriving untagged frames. For an access port the PVID is simply its VLAN.

The distinction that matters: access and trunk ports decide tagging by \emph{role}; a hybrid port decides per VLAN. That is what lets a hybrid port send VLAN 10 untagged to a PC while sending VLAN 20 tagged to a phone on the same cable.
\end{keyidea}

\medskip\noindent\textbf{Part C --- Create VLANs and assign access ports.}\par

\begin{lstlisting}[style=dcnterminal]
system-view
vlan batch 10 20 30

interface GigabitEthernet 0/0/1
 port link-type access
 port default vlan 10
quit

interface GigabitEthernet 0/0/2
 port link-type access
 port default vlan 20
quit
\end{lstlisting}

\medskip\noindent\textbf{Part D --- Carry VLANs between switches.}\par
An access port carries one VLAN. Connecting two switches that share three VLANs would need three cables --- unless the link is a trunk.

\begin{lstlisting}[style=dcnterminal]
interface Eth-Trunk 1
 port link-type trunk
 port trunk allow-pass vlan 10 20 30
quit
\end{lstlisting}

\begin{pitfall}
\cmd{port trunk allow-pass vlan 10 20 30} does not \emph{add} to the allowed list on every VRP version --- on some it replaces it. Always confirm with \cmd{display port vlan} rather than assuming your second command extended the first.

And note that VLAN 1 is allowed by default. Leaving it there is a common audit finding, since it carries management traffic across a link that may not need it.
\end{pitfall}

\medskip\noindent\textbf{Part E --- The native VLAN trap.}\par

\begin{keyidea}
A trunk sends one VLAN untagged: the \term{native VLAN}, which is the trunk port's PVID and defaults to VLAN 1.

If the two ends of a trunk disagree about the native VLAN, traffic leaks silently between them. Switch A sends VLAN 1 untagged; switch B receives an untagged frame and assigns it to \emph{its} native VLAN, say VLAN 99. Two supposedly separate VLANs are now joined, with no error message anywhere.

Set the native VLAN explicitly and identically on both ends, and never use a VLAN that carries user data.
\end{keyidea}

\begin{lstlisting}[style=dcnterminal]
interface Eth-Trunk 1
 port trunk pvid vlan 999
 port trunk allow-pass vlan 10 20 30
quit
\end{lstlisting}

\medskip\noindent\textbf{Part F --- Hybrid ports.}\par

\begin{lstlisting}[style=dcnterminal]
interface GigabitEthernet 0/0/5
 port link-type hybrid
 port hybrid pvid vlan 10
 port hybrid untagged vlan 10
 port hybrid tagged vlan 20
quit
\end{lstlisting}

A PC sending untagged frames lands in VLAN 10 and receives VLAN 10 untagged. A device that understands tags can also use VLAN 20 on the same port.

\medskip\noindent\textbf{Part G --- Verify.}\par

\begin{commandmap}
\begin{tabularx}{\linewidth}{@{}>{\ttfamily}p{0.40\linewidth}X@{}}
\toprule
\rowcolor{MLPaleBlue!92}
\textrm{\bfseries Command} & \bfseries Shows \\
\midrule
display vlan & Every VLAN and its member ports \\
display port vlan & Each port's type, PVID, and allowed VLANs \\
display vlan 10 verbose & Which ports carry VLAN 10, tagged or untagged \\
display mac-address vlan 10 & Which hosts have been learned in that VLAN \\
\bottomrule
\end{tabularx}
\end{commandmap}

Then test: hosts in VLAN 10 must ping each other across both switches, and must fail to reach VLAN 20 even in the same subnet.

\begin{troubleshoot}
\textbf{Same VLAN, same switch, ping fails.} Check both ports are access ports in that VLAN with \cmd{display port vlan}.

\textbf{Same VLAN, different switches, ping fails.} The trunk is not allowing that VLAN, or the VLAN does not exist on the second switch. It must be created on both.

\textbf{Different VLANs can reach each other.} Native VLAN mismatch on the trunk --- the leak described in Part E.

\textbf{Trunk works for some VLANs only.} \cmd{allow-pass} replaced rather than extended the list.
\end{troubleshoot}

\begin{tryit}
Set the native VLAN to 10 on one switch and 20 on the other, then ping between VLANs. It will work --- which is a security failure, not a success.

Capture it in Wireshark and explain which frames carry a tag and which do not. Then fix it and confirm the isolation returns.
\end{tryit}
\end{labproject}

\begin{verifybox}
\begin{itemize}
  \item VLANs created on both switches, not just one.
  \item Trunk explicitly permits every VLAN that must cross it.
  \item Native VLAN set identically on both ends and carries no user data.
  \item Same-VLAN connectivity proven across switches, and cross-VLAN isolation proven to fail.
  \item Hybrid port demonstrated with one VLAN tagged and one untagged.
\end{itemize}
\end{verifybox}

\begin{labdeliverable}
Submit both switch configurations, \cmd{display vlan} and \cmd{display port vlan} output, a connectivity matrix showing which host pairs can reach each other and which cannot, and evidence of the native-VLAN mismatch experiment with an explanation of why it is a security problem.
\end{labdeliverable}
